\documentclass[12pt]{article}


\usepackage[francais]{babel}
\usepackage[babel=true]{csquotes}
\usepackage{amsmath,amsfonts,amssymb,graphicx}
\usepackage{enumerate}
\usepackage{calc}
\usepackage{verbatim}
\usepackage{fancyhdr}
\usepackage[colorlinks=true, linkcolor=blue, urlcolor=blue]{hyperref}
\usepackage{hyperref}
\usepackage{lastpage}
\usepackage{setspace}

\usepackage{color}
\definecolor{bleu}{rgb}{.8,.8,0.91}

\usepackage{geometry}
\geometry{top=30mm, bottom=20mm, left=25mm, right=25mm}

\include{psfig.sty}

\renewcommand{\headrulewidth}{0.6pt}
\renewcommand{\footrulewidth}{0.6pt}
\fancyhead{}
\fancyfoot{}
\fancyhead[LO,LE]{Université Laval \\[-.5ex]
Faculté des sciences et de génie\\[-.5ex]
Département de mathématiques et de statistique
}
\fancyhead[RO,RE]{STT-1900\\Méthodes statistiques\\pour ingénieurs}
%\fancyfoot[LO, LE] {\includegraphics[width=1.5cm]{Figures/UL.jpg}}
%\fancyfoot[RO, RE] {\vspace{-0.5cm} \thepage/\pageref{LastPage}}

\pagestyle{fancy}

\setlength{\parindent}{0mm}
\setlength{\parskip}{0mm}

\renewcommand{\baselinestretch}{1.1}

\begin{document}

\vspace*{1mm}
\centerline{\fcolorbox{black}{bleu}{\bf \large Laboratoire 1}}
\vspace*{3.5mm}

\underline{\bf Objectifs}
\begin{itemize}
\item Se familiariser avec le package {\it Rcmdr} (R-Commander) du logiciel {\it R}.
\item Effectuer quelques calculs impliquant les lois normale, du khi-deux, de Student et de Fisher.
\item Lire un jeu de données.
\item Obtenir des statistiques descriptives et tracer un diagramme en boîte et un histogramme pour un jeu de données.
\item Effectuer des tests d'hypothèses et construire des intervalles de confiance pour la moyenne d'une ou de deux populations normales.
\end{itemize}

\vspace{3.5mm}

\underline{\bf Préparation}
\medskip
\begin{itemize}
\item Visionnez les clips 1 à 8 de la sous-section ``R Commander'' de la section ``Contenu et activités'' du site web du cours.
%\item Relisez les exemples 3.1, 4.1, 4.3, 4.4, 5.3, 5.4, 5.5, 6.2, 7.8 et 7.9 des notes de cours.
\item Installez le logiciel {\it R} et son package {\it Rcmdr}  sur votre ordinateur.\footnote{Un clip et un document PDF disponibles sur le site du cours expliquent la procédure pour installer et charger R et Rcmdr. Le logiciel est installé sur les ordinateurs des laboratoires informatiques du pavillon Pouliot: PLT-2370, PLT-2325.}  Ouvrez R, puis chargez le package Rcmdr.
    %Veuillez noter que contrairement à ce qui était le cas quand le clip 1 a été produit, l'option ``Message translation'' doit être cochée lors de l'installation pour que R s'installe en français (elle était cochée par défaut dans les version précédentes du logiciel).
\end{itemize}

\vspace{3.5mm}

\underline{\bf Exercices}\footnote{Les réponses numériques sont données à la dernière page de ce document.}

\begin{enumerate}

\item Soit $Q$, la quantité  de graisse (en litres) utilisée pour lubrifier les appareils de votre usine pendant une année. On vous dit que $Q$ est une variable aléatoire de loi normale de moyenne 100 litres et d'écart-type 20 litres.
\begin{enumerate}
\item Tracez un graphique de la densité de $Q$. %Sauvegardez ce graphique sous la forme d'un fichier .jpeg. Ouvrez ce fichier à l'aide d'un programme comme {\it Paint} (disponible dans les ``Accessoires'' dans Windows 7) et
    Tentez d'estimer les probabilités ou quantiles demandés aux sous-questions (b) à (e)  "à l'oeil" avant de demander à {\it R} de les calculer exactement.
\item Quelle est la probabilité qu'au plus 110 litres de graisse soient requis pour lubrifier les appareils cette année?
\item Quelle est la probabilité que plus de 125 litres de graisse soient requis pour lubrifier les appareils cette année?
\item Quelle est la probabilité qu'entre 95 et 105 litres de graisse soient requis pour lubrifier les appareils cette année? (Vous devrez calculer 2 valeurs, puis faire la soustraction manuellement. Il est aussi possible d'ajouter un signe $-$ dans la fenêtre {\it Script R} entre les 2 probabilités, puis sélectionner le texte et cliquer sur le bouton {\it Soumettre}.)
\item On vous dit qu'une année sur quatre, il faut plus de $\ell$ litres de graisse pour lubrifier les appareils. Que vaut $\ell$?
\end{enumerate}


\newpage

{\ }

\vspace{3mm}

\item Le fichier Excel {\tt laboratoire1a.xls}\footnote{Ce fichier .xls est disponible sous le même onglet ``Laboratoire 1'' que celui où vous avez trouvé le présent fichier .pdf.} contient les noms des 25 joueurs de poker ayant gagné le plus d'argent à ce jeu entre septembre 2012 et août 2013. On y trouve également leurs gains (en dollars américains) durant cette période.\footnote{Selon le site \href{http://pokerdb.thehendonmob.com/rankings}{The Hendon Mob} en date du 11 septembre 2013.}
\begin{enumerate}
\item Importez le fichier dans {\it R Commander} afin de créer un jeu de données que vous appellerez ``Poker''.
\item Quel est le gain moyen de ces 25 joueurs au cours des douze derniers mois?
\item Quels sont le gain médian et les premier et troisième quartiles de ces gains? Approximativement quelle proportion des observations est inférieure à chacune de ces trois valeurs?
\item Quel est l'écart-type des gains? Comment se compare-t-il à l'écart interquartile?
\item Le joueur ayant le $26^{\mbox{e}}$ plus grand gain pour cette même période est Jason Mercier, avec un gain de  $2\,083\,508\,\$ $. Si on ajoute ce joueur au jeu de données, lesquelles des valeurs calculées en (b) à (d) devraient changer, et dans quelle direction? Vérifiez votre intuition en ajoutant cette observation au jeu de données et en recalculant les statistiques descriptives. (Soit vous cliquez sur le bouton Éditer pour ajouter une ligne, soit vous modifiez le fichier Excel et vous l'importez à nouveau.)
%   LA PROCÉDURE SUIVANTE NE MARCHAIT PAS EN 2016 NI EN 2017, mais elle marche en 2018:
%    \footnote{Pour ajouter une nouvelle observation à un jeu de données en {\it R Commander}, cliquez sur le bouton ``Éditer'', ensuite cliquez sur l'observation en dernière rangée du jeu de données et appuyez sur la touche flèche vers le bas ($\downarrow$) de votre clavier. Une $26^{\mbox{e}}$ rangée est ainsi ajoutée au jeu de données et vous pouvez y entrer la valeur des variables en cliquant sur les cellules correspondantes.}
\end{enumerate}

\vspace{5mm}


\item On considère trois variables aléatoires  indépendantes:  $$Z\sim N(0,1),\qquad U\sim\chi^2_{10},\qquad V\sim\chi^2_{16}.$$
    Évaluez les probabilités ou quantiles ci-dessous, en prenant le temps de les retranscrire d'abord sous une forme facile à évaluer à l'aide de {\it R Commander}.
\begin{enumerate}[(a)]
\item $P(V>16)=?$
\item $P(U+V<20)=?$
\item Soit une nouvelle variable $X=Z/\sqrt{V}$. Que vaut $P(X\le 0,25)$?
\item  Quel est le $75^{\text{e}}$ centile de la distribution de la variable $X$ définie en (c)?
\item  Soit une nouvelle variable $Y=U/V$. Que vaut $P(Y> 0,25)?$
\item  Quelle est la valeur $y$ telle que $P(Y>y)=0,3$, où $Y$ est la variable aléatoire définie en (e)?
\end{enumerate}


\newpage

{\ }
\vspace{3mm}
\begin{comment}
 Ancienne version de la question: inappropriée, car les données ne sont pas indépendantes (série temporelle), la population (juillet 2013) est trop petite, la pige de 16 jours  est invraisemblable.

\item Le fichier {\tt laboratoire1b.xls}  contient deux feuilles de données. Cette question concerne les données de la première feuille ``Labo1, Q4''. On y trouve 16 observations de la température maximale quotidienne observée à la station météorologique de l'aéroport Jean-Lesage de Québec au cours du mois de juillet 2013.
\begin{enumerate}
\item Faites un diagramme en boîte et un histogramme de densité pour ces données.
\item La distribution normale semble-t-elle être une distribution raisonnable pour ces données? Pourquoi?
\item Au seuil de 5\%, y a-t-il de l'évidence que la température maximale quotidienne moyenne en juillet 2013 est supérieure à 25 degrés?
\item Construisez un intervalle de confiance à 90\% pour la température maximale quotidienne moyenne.
\end{enumerate}

Nouveau jeu de données:
set.seed(14632)
x<-round(rnorm(50,mean=4,sd=0.5),2)
y<-round(((0.75*x/pi)^(1/3))*20+rnorm(50,mean=4,sd=0.3),2)
Melon<-1:50
tableau<-data.frame(Melon,x,y)
colnames(tableau)<-c("Melon","Poids","DiametreMax")
write.csv2(tableau,file="question4.csv",quote=F)

\end{comment}

\item Le fichier {\tt laboratoire1b.xls}  contient deux feuilles de données. Cette question concerne les données de la première feuille ``Labo1, Q4''. On y trouve le poids (kg) et le diamètre maximal (cm) de 50 melons d'eau choisis au hasard dans la livraison  d'un fournisseur de Californie.
\begin{enumerate}
\item Faites 3 diagrammes permettant d'étudier la distribution du poids des melons.
\item La distribution normale semble-t-elle être une distribution raisonnable pour le poids des melons? Pourquoi?
\item Au seuil de 5\%, les données nous permettent-elles de conclure que le poids moyen d'un melon est inférieur à 4 kg?
\item Construisez un intervalle de confiance à 90\% pour le poids moyen d'un melon.
\end{enumerate}

\vspace{5mm}

\item Considérons les données de la seconde feuille (Labo1, Q5) du fichier {\tt labo\-ratoire1b.xls}. Il s'agit des résultats d'une expérience menée dans une manufacture où l'on a demandé à 10 soudeurs de tester deux sortes de fil de soudure, disons le fil de type A et le fil de type B. \\

    Chaque soudeur a produit une soudure avec chacun des deux types de fil, ensuite chaque soudure a été soumise à un test de torsion et le couple de torsion menant à la rupture de la soudure a été mesuré. Dans le fichier Excel, la première colonne contient le numéro du soudeur, la seconde le couple de torsion menant à la rupture avec le fil de type A et la troisième le couple de torsion menant à la rupture avec le fil de type B. On peut supposer que les couples de torsion sont des variables aléatoires normalement distribuées. \\

    Au seuil de 5\%, les données présentent-elles des preuves qu'un type de fil mène à des soudures plus solides que l'autre?



\end{enumerate}


%\newpage


\vspace{2mm}

\underline{\hspace{16cm}}
\vspace{2mm}
%%%%%%%%%%%%%%%%%%%%%%%%%%%%%%%%%%%%%%%%%%%%%%%%%%%%%%%%%%%%%%%%%%%%%%%%%%%%%%%%%%

{\bf Réponses}

\small

\begin{enumerate}
\item  \begin{enumerate}
\item[(a)] Cliquez sur Distribution/Distributions continues/Distribution normale/Graphe.
\item[(b)] $P(X<110)=0,6914625$
\item[(c)] $P(X>125)=0,1056498$
\item[(d)] $PX<105)-P(X<95)=0.5987063-0.4012937=0,1974127$
\item[(e)] $P(X>\ell)=0,25\;\Rightarrow\;\ell=113,4898$ litres
\end{enumerate}

\newpage
\vspace*{1mm}
\item  \begin{enumerate}
\item[(a)] Données/Importer des données/Depuis un fichier Excel
\item[(b)] $\overline x=3~344~547~\$$
\item[(c)] $\widetilde x=2~901~462\$;~ Q_1=2~377~821\$;~Q_3=3~444~523\$ $;~~Proportions: $ 50\%;~25\%;~75\%$
\item[(d)] $s$:~1~455~645~\$;~~EIQ:~1~066~702~\$; il s'agit de deux mesures de la dispersion, et dans ce cas-ci l'écart-type est environ 1,36 fois l'EIQ.
\item[(e)] Toutes les quantités calculées ont diminué. On pouvait le deviner pour les mesures de localisation, mais pas vraiment pour les mesures de dispersion.
\end{enumerate}

\item Utilisez le menu Distributions continues.
 \begin{enumerate}[(a)]
 \item $P(V>16)= 0,4529608$
 \item $P(U+V<20)=P(\chi^2_{26}<20)=0,2084435$
\item $P\left(\frac{Z}{\sqrt{V}}\le {0,25}\right)=P\left(\frac{Z}{\sqrt{V/16}}\le \frac{0,25}{\sqrt{1/16}}\right)=P\left(t_{16}\le 1\right)=0,8339025$.
\item  $P\left(X\le q\right)= P\left(\frac{Z}{\sqrt{V/16}}\le \frac{q}{\sqrt{1/16}}\right)= P\left(t_{16}\le \frac{q}{\sqrt{1/16}}\right)=0,75\\[3mm]
    \Rightarrow \frac{q}{\sqrt{1/16}}=0,6901323\Rightarrow q=0,1725331$.
\item  $P\left(\frac{U}{V}> 0,25\right)=P\left(\frac{U/10}{V/16}> \frac{0,25(1/10)}{1/16}\right)=P\left(F_{10;16}> 0,4\right)=0,9274445$.
\item  $P\left(Y> y\right)= P\left(\frac{U/10}{V/16}> \frac{y(1/10)}{1/16}\right)=0,3\Rightarrow 1,6y=1,318637\Rightarrow y=0,8241483$.
\end{enumerate}

\item
 \begin{enumerate}
\item[(a)]  $\mbox{ }$

\vspace{-4mm}

\includegraphics[width=3cm]{hist.png}\hfill\includegraphics[width=3cm]{STT-1900/Diapos/Images/Labos/qq.png}\hfill\includegraphics[width=3cm]{STT-1900/Diapos/Images/Labos/box.png}
\item[(b)] Oui. L'histogramme montre une distribution unimodale raisonnablement symétrique. Le diagramme quantile-quantile est relativement linéaire, aucun point ne sort des bornes pointillées.  Le diagramme en boîte ne révèle aucune donnée aberrante, les moustaches sont approximativement de même longueur et la ligne de la médiane est à peu près au centre de la boîte, ce qui suggère une distribution symétrique.
\item[(c)] [Statistiques/Moyennes/test t univarié]. Non, le seuil observé (valeur P) est $0,1787>5\%$, donc on ne peut pas rejeter $H_0:~\mu=4$ en faveur de $H_1:~\mu<4$ au seuil de 5\%. (On peut aussi dire que $t_{0}=-0,92916 \not< -t_{0,05;49}=-1.676551$, donc $H_0$ non rejetée).
\item[(d)] $[3,81379; 4,05341]$
\end{enumerate}

\item  Il faut effectuer un test t (de Student) pour données appariées (les échantillons A et B ne sont pas indépendants, les deux mesures sont prises par le même soudeur) pour confronter $H_0:~\mu_A=\mu_B$ et $H_1:~\mu_A\neq\mu_B$. Le seuil observé (valeur P) est de $0,01374<0,05$, donc on rejette $H_0$ au seuil de 5\% et on conclut que les deux types de fil donnent des soudures de solidités différentes. (Le fil B est significativement plus solide.)
\end{enumerate}



\end{document}

